\summarychapter{38}{To Earth's Remotest People}{Protestant Missions}

\begin{inbrief}
In 1800 Protestant Christianity hardly existed outside Europe and America; a century later it was in
almost every land. The nineteenth century was the great age of Christian expansion, in sheer scale
without parallel in history. The chapter tells of William Carey, the ``Father of Modern Missions,''
of the pioneers who followed, of David Livingstone in Africa, and of the voluntary society, the new
instrument that made a world mission possible.
\question{What accounts for the sudden burst of Protestant effort to win the whole world for Christ?}
\end{inbrief}

Late in the eighteenth century, in an English village, a sign over a humble workshop read
``Secondhand shoes bought and sold.'' Inside, William Carey mended boots and, when he could, studied
Latin and Greek. Over the bench hung a rough map of the world on which he had written facts from the
voyages of Captain Cook and other travelers. A friend called the shop ``Carey's College.'' Bench and
map stand for the awakening interest in distant peoples. At the time Asia was almost untouched by the
gospel, apart from traces in India and the Dutch East Indies, and Africa, except for the ancient
Copts of Egypt and Ethiopia, was the ``dark continent.'' After eighteen centuries Christianity was far
from a world religion. Kenneth Scott Latourette, the leading historian of its expansion, judges that
never before had any set of ideas, religious or secular, been spread so far by so large a body of full-time workers supported by the free gifts of so many millions.

\ssection{The Pioneer in Modern Missions}
In Protestantism's first century the Catholic powers Spain and Portugal led Europe's expansion, and
its outstanding missionaries were Xavier, Las Casas and Ricci. New lands opened to Protestants only
after the Spanish Armada was defeated (1588) and Britain and the Netherlands rose as colonial powers. The first to go were the Pietists. The Moravians, however, sought out individuals in some
European colony, and the groups they gathered were tiny islands in a sea of heathenism. Carey thought
on a larger scale: the evangelization of whole countries and the results of an entire people turning Christian. In his view the missionary from abroad can do only a little, so that
developing a local ministry is the first of all missionary concerns, and that the faith has to strike root in the customs and culture of the country that receives it.

He was an unlikely candidate for greatness: a cobbler married to a poor, uneducated girl, often
without enough to eat, but with a ravenous hunger for knowledge, who would skip a meal to afford a book and whose heroes were Columbus and Cook. Converted in 1779 through a fellow shoemaker and
baptized as a believer in 1783, he became pastor of the Baptist chapel at Moulton, supporting his
family by teaching and shoemaking. Among the Baptists he met Andrew Fuller, a strong Calvinist who
broke with those who thought earnest evangelism inconsistent with election and who complained that
preachers had sunk into ``a compromising way of dealing with the unconverted.'' From Fuller's
teaching that all are obliged to repent and believe, Carey drew the inference that those entrusted
with the gospel are obliged to take it everywhere. In 1792 he published his \emph{Enquiry}, which argued that Christians are obliged to use practical means to convert the heathen, an epoch-making book. It answered the five objections commonly raised:
the heathen lands' distance, their barbarism, the danger, the difficulty of support and the unknown
languages. None of these stopped the merchants. It needs only, he wrote, as much love for the souls
of our fellow creatures as traders have ``for the profits arising from a few otter skins.'' He ended
with practical proposals. Carey and Fuller read Jesus' command to preach to all the world, and Paul's
word that every knee shall bow, as addressed to themselves. In October 1792 they and eleven
colleagues founded the Baptist Missionary Society, and the Carey family left for India the following year. His wife Dorothy, who had meant to stay behind, yielded to a last appeal, regretted it almost
at once, and spent the rest of her days confined to her room, embittered and possibly insane.

The British East India Company, virtual ruler of India since 1763, cared for profits and thought
sending missionaries the maddest and most dangerous of projects. Refused leave to live in Calcutta,
Carey settled at Serampore under the Danish flag and worked in Bengal as foreman at an indigo works,
a post that occupied three months of the year and left the rest for oriental languages. In 1799
Joshua Marshman and William Ward joined him, and for a quarter of a century the three built a
network of stations in and beyond Bengal. They plunged into Hindu thought, which they considered
essential equipment and no distraction: the preacher must be understood, and the thought-world of a
nation, like its souls and bodies, is to be brought captive to Christ. By 1824 six whole Bibles and portions in twenty-four further languages had been translated under Carey's direction, and he had published grammars, dictionaries
and translations of Eastern books.

\ssection{The Contagion of Missionary Service}
The example of the Serampore trio was contagious. The major churches had thought missions needless
and hopeless; now voices everywhere declared it every Christian's duty to share in the conversion of
the world, for the gospel is not the private possession of Europeans. The pioneers number hundreds:
Henry Martyn and Alexander Duff (India), Robert Morrison (China), John Williams (the South Seas), Adoniram Judson (Burma), Allen Gardiner (Tierra del Fuego), Robert Moffat in South Africa. Scores
more are forgotten because they died within months of fever or at the hands of hostile tribes.

For three decades the effort was almost wholly evangelical, and naturally so. Traditional churchmen
saw the minister's task as nurturing the faith planted at baptism in the whole parish and could not
imagine preaching in a tribal society, while holders of a rigid predestination gave no thought to the
elect in India or China. For Evangelicals to preach was to summon sinners to God by faith in Christ, and they saw no difference in principle between ``baptized heathen'' in Britain and non-Christians
overseas. It was only in the 1820s and 1830s that missions became a normal part of British church life, as evangelical values spread and Britain came to be seen as a Christian nation with responsibilities abroad. Thousands of sermons and hundreds of hymns carried the vision, such as
``From Greenland's Icy Mountains'' by Reginald Heber, who gave his life as bishop of Calcutta.

Two emphases produced this world vision. One was evangelical: people who lack faith in Christ are lost, and Christ has charged each generation to spread the news of salvation. The other was prophetic.
Following Jonathan Edwards, many believed that the whole earth would come to know the Lord, as the prophets had promised, in preparation for Christ's coming reign on earth, a belief called
millennialism. Joined to the optimism of the Age of Progress, the goal seemed attainable, and the
Student Volunteer Movement took as its watchword ``the evangelization of the world in this
generation.'' The vision was kept fresh by reports from the field, none more stirring than those of
David Livingstone (1813--1873) on the spiritual darkness of Africa and the horrors of the Arab slave
trade. He came of hardy Scottish stock, resolved at nineteen to give his life to ``the alleviation
of human misery,'' trained as a doctor, and went out drawn by the fame of Robert Moffat.

\ssection{Livingstone in Africa}
He arrived in 1841 and for ten years did ordinary mission work. But he had an explorer's mind and
was always drawn on by ``the smoke of a thousand villages'' that had never seen a missionary. His
first great journey took him through the jungle to Angola on the Atlantic and then, since he refused to desert his African carriers, back over the whole breadth of Africa to Quilimane on the Indian Ocean. His patience with Africans spared him any resort to force, and he recorded what he saw of science and geography with great precision. With that single journey the interior of Africa was opened to the modern world. Yet he was never only a traveler. ``I place no value on anything I have or may possess,'' he
wrote, ``except in relation to the kingdom of Christ.'' What moved him most was ``this open sore of
the world,'' the slave trade of central Africa. In 1857 he told the students of Cambridge, ``I go
back to Africa to try to make an open path for commerce and Christianity.'' That did not make him a
forerunner of colonial exploiters. He saw that the trade depended on African participation: where
slave raiding was the road to wealth, the temptation to raid weaker neighbors was constant. Only if
Africans could be drawn into legitimate commerce, exchanging the produce of their fields and forests
for the goods of the white man, would the wicked trade stop.

How was the vision turned into action? The denominations were governed by bishops, presbyteries or
congregations, each claiming the Bible, and people had suffered and bled for each form. But none of
them could launch a world mission. Its supporters were driven to another form of cooperation.

\ssection{Creation of the Voluntary Society}
Again Carey was the pioneer. Asking what a trading company would do, his \emph{Enquiry} proposed a
company of serious Christians, laymen and ministers, with a committee to gather information, raise
funds and find suitable missionaries. The voluntary society, whose early form was the missionary society, transformed nineteenth-century Christianity. Devised to meet a need, not for theological
reasons, it nevertheless weakened the existing systems of church government. It allowed action across denominations, since Baptist, Methodist, Anglican and Congregationalist could join forces for a defined purpose without raising questions of church structure. And it shifted
power toward the laity: ordinary men, and later women, held key posts unthinkable elsewhere in the
church. The London Missionary Society declared in 1795 that its design was not to send
``Presbyterianism, Independency, Episcopacy or any other form of church government'' but the gospel;
one founder called for ``the funeral of bigotry.'' Through the societies the enthusiast who collected
a penny a week and passed round the missionary magazine shared fully in the work, and from such
circles candidates came forward. As the American Rufus Anderson wrote in 1834, only in this century
were the evangelical churches really organized for the conversion of the world.

In the United States the first society for work abroad, the American Board of Commissioners for Foreign Missions, dates from 1810 and owed its start to students at the Congregational Andover Theological
Seminary. Their leader, Samuel J.~Mills, had started at Williams College a small society whose
members pledged their lives to missionary service; at Andover it included Adoniram Judson. Students
then, as often later, laid the foundation of an advance. Baptists, Presbyterians and others soon
formed their own agencies, and the conversion of the heathen became a major concern of congregations
everywhere, kept alive by local societies and women's organizations, children's days, visits of
missionaries on furlough, special offerings and finally a place in the regular church budget. By the
end of the century almost every Christian body, from the Russian Orthodox Church to the Salvation Army, and nearly all countries had a share in the enterprise.

The early missionaries were often unaware of the conflicts the gospel caused in societies unlike their own. Most of them took the Western form of the faith for the faith itself, so that to make an Indian or a
Malay a Christian was largely to make him into a Portuguese or a Dutchman. Condemning this is easy,
Shelley says, yet every Christian and every Christian society mixes cultural tradition with the
faith, a problem as old as missions. The results can still be disastrous: if the faith arrives in
Western habits that others find shocking, such as eating meat or freer contact between the sexes, it
is condemned before it is examined. On the other hand the Christian community has a distinctiveness
that comes from the gospel itself, which is a revolutionary power, and to disguise that is to turn
the faith into something else.

\ssection{The Marks of Modern Christianity}
The movement kept expanding, and it brought out several marks of modern Christianity. First, the
expansion came by voluntaryism, not compulsion. Since Constantine the spread of the faith had enjoyed
the support of Christian rulers and often meant mass conversions; these missionaries, with few
exceptions, went without state support or control and advanced only by persuasion, so that mission
and religious freedom went together. Second, it drew on the money and abilities of ordinary Protestants. Earlier missions were led by monks and bishops, whereas the societies were organized on
the broadest base, and the church's primary task, preaching the gospel, fell to its people. Third,
humanitarian work went with the preaching. The agencies founded schools, hospitals and training
centers for nurses and doctors, reduced languages to writing, translated the Bible and other Western
books, and introduced public health measures and better farming. Some of this served conversion;
much sprang simply from need that a Christian conscience could not pass by. By all this the missionary movement put the gospel back at the center and regained something the divisions of the Reformation had hidden, the catholicity of the church: a universality like that
of the map over Carey's bench, reaching out to new peoples in new lands.

\begin{keyfacts}
\keypeople{William Carey; Andrew Fuller; Joshua Marshman; William Ward; Henry Martyn; Robert Morrison;
Adoniram Judson; Alexander Duff; Robert Moffat; Reginald Heber; David Livingstone; Samuel J.~Mills;
Kenneth Scott Latourette.}
\keydates{%
  \kd{1792}{Carey's \emph{Enquiry}; Baptist Missionary Society founded}
  \kd{1793}{Carey sails for India}
  \kd{1795}{London Missionary Society founded}
  \kd{1799}{Marshman and Ward join Carey at Serampore}
  \kd{1810}{American Board of Commissioners for Foreign Missions founded}
  \kd{1841}{David Livingstone arrives in Africa}
  \kd{1857}{Livingstone's appeal at Cambridge for ``commerce and Christianity''}
  \kd{1873}{Death of Livingstone}}
\keyterms{Modern missions; local ministry; Serampore trio; millennialism; ``the evangelization of the
world in this generation''; voluntary society; missionary society; lay leadership; voluntaryism;
humanitarian ministries; catholicity.}
\end{keyfacts}
